\scanpage{81}
% printed form
% letterhead
UNIVERSITE\\
DES SCIENCES ET TECHNIQUES\hfill Montpellier, le 1er décembre 1977\\
\uline{DU LANGUEDOC}\\
JP/MD\\
Diffusion : Liste A\\
n° 6827

AVIS DE SOUTENANCE D'UNE THESE DE DOCTORAT DE SPECIALITE

SCIENCES BIOLOGIQUES - Mention PHYSIOLOGIE ANIMALE

\uline{Nom et prénom du candidat} : FAVRE Daniel

\uline{Titre de la thèse} : Ontogenèse de l'innervation des récepteurs vestibulaires chez le chat. Etude neuro-anatomique et ultrastructurale.

\uline{Date, heure et lieu de soutenance prévus} : 10 décembre 1977 à 11 heures, salle de cours 5

\uline{Nom du Directeur de thèse} : Monsieur A. SANS

\uline{Désignation du laboratoire où la thèse est préparée} : Laboratoire de neurophysiologie sensorielle.

\uline{Analyse de l'ouvrage} : Etude de l'ontogenèse de l'innervation afférente et efférente des cellules sensorielles vestibulaires. Mise en évidence, au cours de la maturation, d'une modification structurale des contacts synaptiques entre les calices nerveux et les cellules de type I. Discussion relative à l'origine phylogénétique des cellules labyrinthiques chez les mammifères.

\uline{Durée approximative des recherches ayant abouti à la thèse} : 2 ans

\uline{Composition du jury proposé} : M. R. MARTY - Président\\
M. J. LEGRAND\\
M. R. PUJOL\\
M. A. SANS

Le Secrétaire Général,\\
\handwritten{\unsure{E. Siau|E. Si}}\\
\uline{E. SIAU}

% archive sheet number: 80
